
\section{Demostració}
\iniciTest{t2}{d,d,c,d,b,a,d,b,a,b,d}

\begin{enumerate}
\pregunta Considereu l'enunciat \enquote{ per a tots els nombres enters
$a$ i $b$, si $a+b$ és parell, aleshores $a$ i $b$ són
parells}. Llavors:

\begin{enumerate}
\opcio{a}{ L'enunciat recíproc és \enquote{ per a tots els nombres
enters $a$ i $b$, si $a$ o $b$ no és parell, aleshores $a+b$ no és
parell}.

}
\opcio{b}{ L'enunciat contrarecíproc és \enquote{ per a tots els
nombres enters $a$ i $b$, si $a$ i $b$ són parells, aleshores $a+b$ és
parell}.

}
\opcio{c}{ El contrarecíproc és fals i un contraexemple és, per
exemple, $a=4$ i $b=7$.

}
\opcio{d}{ L'enunciat contrari és \enquote{ Hi ha nombres $a$ i $b$
tals que $a+b$ és parell però $a$ i $b$ no són
parells}\ i és vertader. 
}
\end{enumerate}
\feedback

\pregunta Suposem que $a$ i $b$ són nombres reals. Volem provar que si $0<a<b
$ aleshores $a^{2}<b^{2}$. Llavors, quina de les següents afirmacions
és falsa:

\begin{enumerate}
\opcio{a}{ La prova directa consisteix a prendre com a hipòtesi $0<a<b$ i com a
tesi $a^{2}<b^{2}$.

}
\opcio{b}{ La prova cap enrere consisteix a prendre com a punt de partida la
conclusió $a^{2}<b^{2}$ i arribar a deduir que $a<b$ suposant que $a$ i
$b$ són positius, i, després, refer la prova procedint de forma directa.

}
\opcio{c}{ La prova per contrarecíproc consisteix a prendre com a hipòtesi
$a^{2}\geq b^{2}$ i com a tesi que no es compleix $0<a<b$.

}
\opcio{d}{ La prova per reducció a l'absurd consisteix a prendre com a hipòtesi
$a\geq b>0$ i $a^{2}<b^{2}$ i com a tesi una contradicció. 
}
\end{enumerate}
\feedback

\pregunta Sigui $x=\dfrac{y}{y^{2}+1}$. Llavors
\[
y-x=y-\dfrac{y}{y^{2}+1}=\dfrac{y^{3}}{y^{2}+1}=\dfrac{y}{y^{2}+1}\cdot
y^{2}=xy^{2}\text{.}%
\]
i, per tant, l'enunciat: $\exists x\in\mathbb{R~}\forall y\in\mathbb{R~}%
xy^{2}=y-x$ és un teorema.

\begin{enumerate}
\opcio{a}{ La prova del teorema és correcta.

}
\opcio{b}{ La prova correspon al teorema: $\forall y\in\mathbb{R}$ $\exists
x\in\mathbb{R~}x^{2}y=y-x$

}
\opcio{c}{ La prova del teorema és incorrecta perquè no es pot definir $x$
abans de $y$. 

}
\opcio{d}{ Cap de les respostes és vertadera.
}
\end{enumerate}
\feedback

\pregunta L'enunciat: \enquote{ Suposem que $m$ és un nombre enter
parell i $n$ és un nombre enter senar. Aleshores $n^{2}-m^{2}%
=n+m$}\ és un teorema?

\begin{enumerate}
\opcio{a}{ \enquote{$n^{2}-m^{2}=n+m$}\ és la
condició necessària.

}
\opcio{b}{ \enquote{$m$ és un nombre enter parell i $n$ és un
nombre enter senar}\ és condició necessària i suficient.

}
\opcio{c}{ És teorema perquè $m=2k$ i $n=2k+1$, on $k\in\mathbb{Z}$, i
$n^{2}-m^{2}=\left(  2k+1\right)  +2k=m+n$.

}
\opcio{d}{ No és teorema i, $m=2$ i $n=3$ és un contraexemple. 
}
\end{enumerate}
\feedback

\pregunta Quina de les següents respostes és falsa?

\begin{enumerate}
\opcio{a}{ Per demostrar que $\forall x\in\mathbb{R}\left(  \exists y\in
\mathbb{R}(x+y=xy)\leftrightarrow x\neq1\right)  $ fem la prova següent:
Considerem un nombre real arbitrari $x$, aleshores existeix un nombre real
$y=\dfrac{x}{x-1}$ si i només si $x\neq1$ i també es compleix%
\[
x+y=x+\dfrac{x}{x-1}=\frac{x^{2}}{x-1}=x\cdot\frac{x}{x-1}=xy\text{.}%
\]


}
\opcio{b}{ Per demostrar que $\forall x\in\mathbb{R}\left(  \exists y\in
\mathbb{R}(x+y=xy)\leftrightarrow y\neq1\right)  $ fem la prova següent:
Considerem $x=\dfrac{y}{y-1}$ que existeix si i només si $y\neq1$.
Aleshores es compleix%
\[
x+y=\dfrac{x}{x-1}+y=\frac{xy}{x-1}=\frac{x}{x-1}\cdot y=xy\text{. }%
\]


}
\opcio{c}{ Per provar que $\exists z\in\mathbb{R~}\forall x\in\mathbb{R}^{+}\left(
\exists y\in\mathbb{R~}y-x=y/x\longleftrightarrow x\neq z\right)  $ fem:
Suposem que existeix un nombre real $a$ tal que
\[
y-a=\frac{y}{a}\Longrightarrow y=\dfrac{a}{\dfrac{a-1}{a}}%
\]
aleshores $y$ existirà si i només si $a\neq0$ i $a\neq1$. D'aquí
s'obté que per a qualsevol $x\in\mathbb{R}^{+}$ existeix $y\in\mathbb{R}$
si i només si existeix $z=1$ i $x\neq1$.

}
\opcio{d}{ Per demostrar que per a cada nombre enter $n$, $6~|~n$ si i només si $2~|~n$ i
$3~|~n$ fem la prova següent: Sigui $n$ qualsevol nombre enter. La
condició $2~|~n$ i $3~|~n$ és necessària perquè $n=2p$ i
$n=3q$ i $p,q\in\mathbb{Z}$. Llavors $n=3\left(  2p\right)  -2\left(
3q\right)  =6\left(  p-q\right)  $ i $6~|~n$. La condició $6~|~n$ és
suficient perquè $n=6k=3(2k)$, $k\in\mathbb{Z}$, i, per tant, $3~|~n$, i
$n=2\left(  3k\right)  $ i així $2~|~n$.
}
\end{enumerate}
\feedback

\pregunta Quina de les següents respostes és correcta?

\begin{enumerate}
\opcio{a}{ Sabem que hi ha nombres primers; per exemple, $2$ és primer.
Suposem ara que només hi ha un nombre finit de nombres primers
$p_{1},p_{2},...,p_{n}$, $n\in\mathbb{Z}^{+}$. Llavors $q=p_{1}p_{2}\cdots
p_{n}+1$ té algun divisor primer. Si tots els primers fossin exactament
$p_{1},p_{2},\ldots,p_{n}$, existiria un $k$, amb $1\leq k\leq n$, tal que
$p_{k}$ dividiria $q$, o sigui%
\[
q=p_{k}r
\]
D'aquí%
\begin{align*}
p_{1}p_{2}\cdots p_{n}+1 &  =p_{k}r\\
p_{k}\left(  r-p_{1}\cdots p_{k-1}p_{k+1}\cdots p_{n}\right)   &  =1
\end{align*}
i, per tant, $p_{k}$ divideix $1$, però això és una
contradicció perquè $p_{k}$ és primer. Per consegüent, ha
d'existir algun primer que no és a la llista inicial. Per tant, hem
demostrat per reducció a l'absurd que hi ha infinits nombres primers. 

}
\opcio{b}{ Sigui $n$ qualsevol nombre enter. Si $n$ és senar, aleshores
$n=2k+1$, $k\in\mathbb{Z}$. Llavors,%
\[
n^{2}=\left(  2k+1\right)  ^{2}=2\left(  2k^{2}+2k\right)  +1
\]
i, per tant, $n^{2}$ és senar perquè $2k^{2}+2k\in\mathbb{Z}$.
D'aquesta manera hem demostrat per contrarecíproc que per qualsevol enter
$n$, si $n$ és parell aleshores $n^{2}$ també ho és.

}
\opcio{c}{ Considerem dos nombres reals qualssevol $x$ i $y$. Suposem que $xy$
és irracional i $x$ és racional. Llavors $x=p/q$, on $p,q\in
\mathbb{Z}$ i $q\neq0$. Suposem ara també que $y$ és racional, llavors
$y=r/s$, on $r,s\in\mathbb{Z}$ i $s\neq0$. Per tant, es té $xy=pr/qs$ i
$qs\neq0$. Aleshores $xy$ és racional i això és una
contradicció. Per consegüent, $y$ és irracional. D'aquesta manera
hem demostrat per contrarecíproc que si $x$ o $y$ és racional,
aleshores $xy$ és racional.

}
\opcio{d}{ Sigui $n$ un nombre enter senar qualsevol. Aleshores $n=2s+1$,
$s\in\mathbb{Z}$. Suposem que $s$ és parell, aleshores $s=2p$,
$p\in\mathbb{Z}$. Llavors $n^{2}=\left(  2s+1\right)  ^{2}=\left(
4p+1\right)  ^{2}=8\left(  2p^{2}+p\right)  +1$ i, per tant, $n^{2}=8k+1$,
$k=2p^{2}+p\in\mathbb{Z}$. D'aquesta manera hem demostrat per casos que si $n$
és un nombre enter senar qualsevol, aleshores existeix un enter $k$ tal
que $n^{2}=2k+1$.
}
\end{enumerate}
\feedback

\pregunta Quina de les següents respostes és falsa?

\begin{enumerate}
\opcio{a}{ Considerem $x$ un nombre real arbitrari, i suposem que $x\neq2$. Ara
considerem $y=\dfrac{x}{2-x}$, que existeix ja que $x\neq2$. Aleshores es
té%
\[
\frac{2y}{y+1}=\frac{\dfrac{2x}{2-x}}{\dfrac{x}{2-x}+1}=\frac{\dfrac{2x}{2-x}%
}{\dfrac{2}{2-x}}=\frac{2x}{x}=x\text{.}%
\]
Per veure que aquesta solució és única, suposem que existeix $z$
tal que $2z/(z+1)=x$. Aleshores $2z=x\left(  z+1\right)  $, i d'aquí surt
$z\left(  2-x\right)  =x$. Com que $x\neq2$ podem dividir els dos costats per
$2-x$ per obtenir $z=x/\left(  2-x\right)  =y$. D'aquesta manera hem demostrat
que per a cada nombre real $x$, si $x\neq2$, hi ha un nombre real únic $y$
tal que $2y/\left(  y+1\right)  =x$.

}
\opcio{b}{ Sabem que $\sqrt{2}$ és irracional. Aleshores $\left(  \sqrt
{2}\right)  ^{\sqrt{2}}$ és racional o bé irracional. Si $\left(
\sqrt{2}\right)  ^{\sqrt{2}}$ és racional, prenem $a=b=\sqrt{2}$ es té
que hi ha dos irracionals tal que $a^{b}=\left(  \sqrt{2}\right)  ^{\sqrt{2}}$
és racional. Per altra banda, si $\left(  \sqrt{2}\right)  ^{\sqrt{2}}$
és irracional, prenem $a=\left(  \sqrt{2}\right)  ^{\sqrt{2}}$ i
$b=\sqrt{2}$ es té que hi ha dos irracionals tals que $a^{b}=\left(
\left(  \sqrt{2}\right)  ^{\sqrt{2}}\right)  ^{\sqrt{2}}=\left(  \sqrt
{2}\right)  ^{2}=2$ és racional. D'aquesta manera hem demostrat que hi ha
nombres irracionals $a$ i $b$ tals que $a^{b}$ és racional.

}
\opcio{c}{ Suposem que $mn$ és múltiple de 3 i que $n$ no és
múltiple de 3. Llavors, $mn=3k$, $k\in\mathbb{Z}$, i d'aquí, surt que
$m$ necessàriament és múltiple de 3 perquè $n$ no ho és.
Així hem provat que $m$ o $n$ és múltiple de 3 és
condició necessària perquè $mn$ és múltiple de 3. Suposem
ara que $m$ és múltiple de 3; es prova anàlogament si $n$ és
múltiple de 3. Llavors, $m=3p$, $p\in\mathbb{Z}$, i d'aquí surt que
$mn=3pn$ i, per tant, $mn$ és mútiple de 3 perquè $pn\in
\mathbb{Z}$. Això prova que $m$ o $n$ és múltiple de 3 és
condició suficient. Per tant, $mn$ és múltiple de 3 si i només si $m$ o $n$
és múltiple de 3 és un teorema.

}
\opcio{d}{ Sigui $n$ el nombre enter més gran. Aleshores, com que $1$ és un
nombre enter, és clar que $1\leq n$. D'altra banda, com que $n^{2}$
també és un nombre enter també es compleix $n^{2}\leq n$ i
d'aquí s'obté que $n\leq1$. D'aquesta manera, com que es compleix
$n\leq1$ i $n\geq1$, aleshores es té $n=1$, i, per tant, $1$ és
l'enter més gran. 
}
\end{enumerate}
\feedback

\pregunta Examina la següent fal·làcia:

(i) Considerem l'equació $\frac{x+5}{x-7}-5=\frac{4x-40}{13-x}$, suposant
que $x\neq7$ i $x\neq13$.

(ii) Operant en el terme de l'esquerra, es pot comprovar que $\frac{x+5}%
{x-7}-5=\frac{4x-40}{7-x}$.

(iii) De (i) i (ii) es dedueix que $\frac{4x-40}{13-x}=\frac{4x-40}{7-x}$.

(iv) Atès que els numeradors són iguals, els denominadors també ho
han de ser. És a dir, de (iii) es dedueix que $7-x=13-x$.

(v) De (iv) es dedueix que 7 = 13. Absurd!

En algun dels passos (i)--(v) hi ha d'haver un error. Quin és ? Per
què ?

\begin{enumerate}
\opcio{a}{ pas (ii) perquè es dedueix $\frac{x+5}{x-7}-5=\frac{40-4x}{x-7}.$

}
\opcio{b}{ pas (iv) perquè es dedueix $\left(  4x-40\right)  \left(
7-x\right)  =\left(  4x-40\right)  \left(  13-x\right)  $ i, d'aquí no
s'obté $7-x=13-x$ llevat que $x\neq10$ 

}
\opcio{c}{ pas (iv) perquè es dedueix $13-x=x-7$

}
\opcio{d}{ pas (v) perquè \ es dedueix $2x=20$ i, per tant, $x=10$.
}
\end{enumerate}
\feedback

\pregunta Una fal·làcia: En qualsevol bossa de bales, totes
les bales són del mateix color.

Demostració per inducció: Sigui $n$ el nombre de bales de la bossa. Si
$n=1$, és evidentment cert. Suposem que és cert per a totes les bosses
de $n$ bales, i considerem una bossa de $n+1$ bales. N'apartem una, i
així tenim una bossa de $n$ bales, que per la hipòtesi d'inducció
seran totes del mateix color. Ens falta provar que la bala apartada també
és del mateix color. L'afegim a la bossa de $n$ bales que havíem
format, i n'apartem una altra. Tornem a tenir una bossa de $n$ bales, que per
hipòtesi d'inducció seran totes del mateix color, i en particular la
darrera bala serà del mateix color que les altres en la bossa. Així
doncs, totes les $n+1$ bales són del mateix color.

\begin{enumerate}
\opcio{a}{ L'error apareix en el pas de $n=1$ a $n=2$: les dues bosses d'una
sola bala que es comparen no tenen cap bala en comú, i per tant no hi ha
cap element compartit que permeti relacionar-ne els colors.

}
\opcio{b}{ No podem fer una prova per inducció perquè la propietat no
està relacionada amb objectes matemàtics.

}
\opcio{c}{ No podem aplicar la hipòtesi d'inducció en el segon cas
perquè $n$ no és qualsevol sinó el que teníem al principi. 

}
\opcio{d}{ Cap de les respostes anteriors és correcta.
}
\end{enumerate}
\feedback

\pregunta Per a tots els $n\geq3$, si $n$ punts diferents d'una circumferència
estan connectats de manera consecutiva amb rectes, llavors els angles
interiors del polígon resultant sumen $\left(  n-2\right)  180%
%TCIMACRO{\U{ba}}%
%BeginExpansion
{{}^o}%
%EndExpansion
$. Quin dels passos següents en la prova per inducció hi ha error?

\begin{enumerate}
\opcio{a}{ Cas base: Suposem que $n=3$. Aleshores el polígon és un
triangle i es compleix perquè se sap que els angles interiors d'un
triangle sumen $180%
%TCIMACRO{\U{ba}}%
%BeginExpansion
{{}^o}%
%EndExpansion
$.

}
\opcio{b}{ Hipòtesi d'inducció: Donats $n$ punts diferents d'una
circumferència estan connectats de manera consecutiva amb rectes, llavors
els angles interiors del polígon resultant sumen $\left(  n-2\right)  180%
%TCIMACRO{\U{ba}}%
%BeginExpansion
{{}^o}%
%EndExpansion
$. 

}
\opcio{c}{ Considerem ara el polígon P format per la connexió de $n+1$
punts diferents $A_{1},A_{2},...,A_{n+1}$ en un cercle, com podeu veure a la
figura 1. Si ens saltem l'últim punt $A_{n+1}$, llavors obtenim un
polígon P amb només $n$ vèrtexs, i per la hipòtesi
d'inducció, els angles interiors d'aquest polígon sumen $\left(
n-2\right)  180%
%TCIMACRO{\U{ba}}%
%BeginExpansion
{{}^o}%
%EndExpansion
$.%
\FRAME{dtbpFU}{7.2664cm}{6.379cm}{0pt}{\Qcb{Figura 1}}{}{logi3.jpg}{\special{ language "Scientific Word";  type "GRAPHIC";
maintain-aspect-ratio TRUE;  display "USEDEF";  valid_file "F";
width 7.2664cm;  height 6.379cm;  depth 0pt;  original-width 7.1939cm;
original-height 6.3087cm;  cropleft "0.009185";  croptop "1";
cropright "1.009185";  cropbottom "0";
filename 'img/logi3.jpg';file-properties "XNPEU";}}


}
\opcio{d}{ Però la suma dels angles interiors de P és igual a la suma dels
angles interiors de P més la suma dels angles interiors del triangle
$A_{1}A_{n}A_{n+1}$. Com que la suma dels angles interiors del triangle és
$180%
%TCIMACRO{\U{ba}}%
%BeginExpansion
{{}^o}%
%EndExpansion
$, podem concloure que la suma dels angles interiors de P és $(n-2)180%
%TCIMACRO{\U{ba}}%
%BeginExpansion
{{}^o}%
%EndExpansion
+180%
%TCIMACRO{\U{ba}}%
%BeginExpansion
{{}^o}%
%EndExpansion
=((n+1)-2)180%
%TCIMACRO{\U{ba}}%
%BeginExpansion
{{}^o}%
%EndExpansion
$.
}
\end{enumerate}
\feedback

\pregunta Per a qualsevol nombre enter positiu $n$, una graella quadrada de
$2^{n}\times\ 2^{n}$ amb qualsevol quadrat eliminat es pot cobrir amb rajoles
en forma de L, com
\[%
\begin{tabular}
[c]{|l|l}\cline{1-1}
& \\\hline
& \multicolumn{1}{|l|}{}\\\hline
\end{tabular}
\]
La figura 2 mostra un exemple per al cas $n=2$. En aquest cas $2^{n}=4$, i per
tant tenim una graella de $4\times4$ i el quadrat que s'ha eliminat està
ombrejat. Les línies pesades mostren com es poden cobrir els quadrats
restants amb cinc rajoles en forma de L.%
\FRAME{dtbpFU}{12.0287cm}{4.145cm}{0pt}{\Qcb{Figura 2}}{}{logi4.jpg}{\special{ language "Scientific Word";  type "GRAPHIC";maintain-aspect-ratio TRUE;  display "USEDEF";  valid_file "F";
width 12.0287cm;  height 4.145cm;  depth 0pt;  original-width 6.8029cm;
original-height 2.5898cm;  cropleft "0";  croptop "1";  cropright "1";
cropbottom "0";  filename 'img/logi4.jpg';file-properties "XNPEU";}}

Quin dels passos següents en la prova per inducció hi ha error?

\begin{enumerate}
\opcio{a}{ Cas base: Suposem que $n=1$. Aleshores la quadrícula és una
quadrícula de $2\times2$ amb un quadrat eliminat, que es pot cobrir
clarament amb una rajola en forma de L.

}
\opcio{b}{ Hipòtesi d'inducció: Sigui $n$ un nombre enter positiu
arbitrari, i suposem que la graella $2^{n}\times2^{n}$ amb qualsevol quadrat
eliminat es pot cobrir amb rajoles en forma de L.

}
\opcio{c}{ Considerem ara una graella $2^{n+1}\times2^{n+1}$ amb un quadrat
eliminat, com es veu a la figura 3. Tallem la graella per la meitat tant
verticalment com horitzontalment, dividint-la en quatre quadrícules
$2^{n}\times2^{n}$.
\FRAME{dtbpFU}{5.6277cm}{5.6014cm}{0pt}{\Qcb{Figura 3}}{}{logi5.jpg}{\special{ language "Scientific Word";  type "GRAPHIC";
maintain-aspect-ratio TRUE;  display "USEDEF";  valid_file "F";
width 5.6277cm;  height 5.6014cm;  depth 0pt;  original-width 3.2159cm;
original-height 3.2159cm;  cropleft "0";  croptop "1";  cropright "1";
cropbottom "0";  filename 'img/logi5.jpg';file-properties "XNPEU";}}

El quadrat que s'ha eliminat està dins d'una d'aquestes quadrícules
i, per tant, per la hipòtesi d'inducció la resta d'aquesta
quadrícula es pot cobrir amb rajoles en forma de L.

}
\opcio{d}{ Les altres quadrícules també es poden cobrir amb rajoles en
forma de L perquè són del tipus $2^{n}\times2^{n}$ i podem aplicar la
hipòtesi d'inducció. 
}
\end{enumerate}
\feedback
\end{enumerate}